\begin{afterword}
\summaryhead

The post wonders whether the Pythagoreans, who kept their teachings secret, had the right idea. The author has argued that science is public by nature, but imagines another future in which science is hidden and guarded by robed cults, so that more people will study it.

The reason is the scarcity effect: what seems hard to get is valued more, and the author says this holds especially for information. People, the author thinks, assume that freely available information is unimportant, so they ignore science and join cults that keep their truths secret. Hidden science would draw people in, satisfy them when they got it, and show up rival cults; it might also be more fun. To the objection that science is already hard to learn, the post answers that the public thinks it is freely available.

The author then disclaims the proposal and asks readers to imagine the robes, so as not to undervalue science, while trusting no robed group that does not show its data. Science, the post ends, is not really free, though people think it is; it should be the other way around.

\responsehead

The post is partly a joke and is judged here as one: a thought experiment with a disclaimer and a practical point. The jokes are plainly flagged. The one serious claim is the explanation of why people ignore science, and that is the part the post supports least.

The claim is that people take freely available information to be unimportant. It is offered with ``I think'' and no evidence of its own. Its support is the scarcity effect of the previous day's post, where the studies of information show forbidden or exclusive information gaining influence: a banned speech moved opinion, and buyers ordered more beef on exclusive news. None of them measures whether people try harder to obtain information. The claim the post needs is about the effort to learn, and that was not measured.

The post raises an objection itself: science is hard to learn and has its own initiations, so why does the scarcity effect not apply already? Its answer, that scarcity works on appearances and the public thinks science is freely available, fits the effect. But it treats difficulty only as a possible source of scarcity. Difficulty is also, by itself, a reason not to study something. At the end the post adds that courses and textbooks are expensive, and calls the situation ``the worst of both worlds.'' So the post grants two plain reasons why few people study science, difficulty and cost, and gives no evidence that favors its own explanation over them.

Other parts hold up better than the post shows. That a costly initiation makes people value what they receive has experimental support: Aronson and Mills (1959) found that a severe initiation made a group seem more attractive to those who went through it. The historical premise, a secret Pythagorean mathematics, rests partly on later tradition; secrecy is attested for some Pythagorean doctrines, not for all. The practical advice keeps the standard of evidence, telling readers to trust no group in robes ``unless they also show you the data.'' Whether imagining the robes changes how anyone values science, the post does not test.

In short: an entertaining thought experiment whose one serious claim, that people ignore science because it seems free, is hedged with ``I think,'' rests on studies that measured the influence of forbidden information rather than the effort to obtain it, and is not weighed against the plainer reasons the post itself grants, difficulty and cost.
\end{afterword}
